\documentclass[aspectratio=169]{beamer}

\usepackage[utf8]{inputenc}
\usepackage[T1]{fontenc}
\usepackage[brazil]{babel}
\usepackage{graphicx}
\usepackage{booktabs}

\usetheme{Singapore}


\title[APS e Planejamento de Visitas]{Aten\c{c}\~ao Prim\'aria \`a Sa\'ude\\}
\subtitle{INF99003 -- Projeto 2 -- Grupo D}
\author{Fabio, Tobias e Vítor}
\date{Semana 1 - 18/09/2026}

\setbeamertemplate{navigation symbols}{}
\setbeamercovered{transparent}
\setbeamertemplate{footline}[frame number]

\begin{document}

% ============================================================
\begin{frame}
  \titlepage
\end{frame}

% ============================================================
\section{Contexto e Problema}

\begin{frame}{Atenção primaria a saude}
  \begin{itemize}
    \item A APS é \textbf{coordenar redes de atenção} para reduzir mortes e internações evitáveis, especialmente doenças crônicas. \\
    \item No Brasil, sua organização ocorre principalmente por meio de \textbf{equipes multiprofissionais} responsáveis por um território definido.
  \end{itemize}
\end{frame}


% ============================================================
\section{Embasamento Te\'orico - Artigos}

\begin{frame}{Evolução da estrutura e resultados da Atenção Primária à Saúde no Brasil entre 2008 e 2019}
  \begin{itemize}
    \item Entre 2008 e 2019, a \textbf{despesa mediana} em APS por habitante coberto aumentou 12,1\%.
    \item A \textbf{cobertura mediana} passou de 98,8\% para 100\%.
    \item A mortalidade por CSAP aumentou 0,2\%.
    \item As \textbf{internações por CSAP} diminuíram 44,9\%.
  \end{itemize}
\end{frame}

% ============================================================
\begin{frame}{Atenção primária à saúde}
  \begin{itemize}
    \item A APS é uma \textbf{estratégia de organização da atenção}.
    \item A APS seletiva oferece ações \textbf{focalizadas} e de baixo custo.
    \item A APS universal e integral articula \textbf{promoção, prevenção, tratamento e acompanhamento}.
    \item A \textbf{Estratégia Saúde da Família} é a principal forma de configuração da Atenção Básica no Brasil.
  \end{itemize}
\end{frame}

% ============================================================
\begin{frame}{A territorialização da Atenção Básica à Saúde do Sistema Único de Saúde do Brasil}
  \begin{itemize}
    \item O território é um \textbf{espaço social, político e dinâmico}.
    \item Cada domicílio deve ser compreendido dentro da \textbf{área de responsabilidade da equipe}.
    \item A municipalização foi essencial, mas \textbf{não resolve sozinha} a organização da rede.
    \item O planejamento de visitas deve considerar o \textbf{território}, e não apenas a localização.
  \end{itemize}
\end{frame}

% ============================================================
\begin{frame}{A atenção primária à saúde na  coordenação  das  redes  de  atenção:  uma  revisão integrativa}
  \begin{itemize}
    \item Uma rede coordenada pela APS pode melhorar a \textbf{qualidade clínica}.
    \item Uma APS fortalecida pode \textbf{racionalizar recursos} e reduzir custos.
    \item O cuidado deve se adaptar às condições \textbf{sociais, demográficas e epidemiológicas} da população.
    \item A \textbf{coordenação efetiva} exige sistemas logísticos e de apoio.
  \end{itemize}
\end{frame}

% ============================================================
\begin{frame}{A Atenção Primária à Saúde no SUS: avanços e ameaças.}
  \begin{itemize}
    \item A APS é responsável por \textbf{coordenar o cuidado}.
    \item A APS e a ESF \textbf{ampliaram o acesso} da população brasileira.
    \item Também contribuíram para melhorar indicadores de saúde e \textbf{reorganizar o cuidado}.
    \item O texto critica propostas de \textbf{financiamento por vouchers}.
  \end{itemize}
\end{frame}

% ============================================================
\begin{frame}{Inovações na atenção Primária em Saúde: o uso de ferramentas de tecnologia de comunicação e informação para apoio à gestão local}
  \begin{itemize}
    \item A Rede OTICS-RIO apoia a \textbf{gestão local} por meio de tecnologias de informação e comunicação.
    \item Reunia \textbf{16 observatórios} e apoiava \textbf{193 unidades} de Atenção Primária.
    \item Os blogs registravam o trabalho das equipes e divulgavam ações de saúde.
    \item A experiência utilizava \textbf{tecnologias simples e de baixo custo}.
  \end{itemize}
\end{frame}

% ============================================================
\begin{frame}{O uso de Tecnologias de Informação e Comunicação em Saúde 
na Atenção Primária à Saúde no Brasil, de 2014 a 2018}
  \begin{itemize}
    \item O estudo analisou o uso de TICS pelas equipes de APS entre 2014 e 2018.
    \item As TICS \textbf{produzem, armazenam, transmitem e protegem informações}.
    \item O uso da \textbf{Telessaúde} aumentou de 32,7\% para 54,6\% das equipes.
    \item Persistiram \textbf{desigualdades regionais} e limitações de infraestrutura.
  \end{itemize}
\end{frame}


% ============================================================
\section{Conclus\~ao}

\begin{frame}{S\'intese}
  \begin{itemize}
    \item A APS territorializada \textbf{é pilar do SUS}
    \item \textbf{Visitas domiciliares} são essenciais para continuidade e equidade
    \item \textbf{Planejamento inteligente} combina:
      \begin{itemize}
        \item \textbf{Otimiza\c{c}\~ao log\'istica}
        \item \textbf{Prioridade cl\'inica}
        \item \textbf{Tecnologia como suporte}
      \end{itemize}
  \end{itemize}
\end{frame}

\begin{frame}{Pr\'oximos Passos}
  \begin{itemize}
    \item \textbf{Semana atual}: revis\~ao bibliogr\'afica e consolida\c{c}\~ao do entendimento do problema.
    \item \textbf{Pr\'oxima etapa}: defini\c{c}\~ao dos requisitos do artefato e desenvolvimento de um \textbf{prot\'otipo inicial}.
  \end{itemize}
\end{frame}

% ============================================================
\begin{frame}{Refer\^encias}
  \begin{columns}[T]
    \column{0.48\textwidth}
    \begin{itemize}
      \item Barros et al. (2022)
      \item Matta \& Morosini (Fiocruz)
      \item Faria (2020)
      \item Rodrigues et al. (2014)
    \end{itemize}
    \column{0.48\textwidth}
    \begin{itemize}
      \item CONASS (2021)
      \item Pinto \& Rocha (2015)
      \item Bender et al. (2024)
    \end{itemize}
  \end{columns}
\end{frame}

% ============================================================
\begin{frame}{}
  \centering
  \vspace{1cm}
  {\Huge Obrigado!}\\[1em]
  {\large D\'uvidas e discuss\~ao}
\end{frame}

\end{document}
